%! Unidad 2 — Disoluciones: ejercicios resueltos y ejercicios finales
%! Fuente: Fuentes/Quimica/Unidad 2 («Ejercicios resueltos», «Practica y calcula», «Piensa y resuelve»).
%! Enunciados adaptados con datos cambiados; soluciones recalculadas. Dificultad: \(\star\) baja, \(\star\star\) media, \(\star\star\star\) alta.
%! En las propiedades coligativas se han cambiado los solutos iónicos por moleculares (salvo en los ejercicios con factor de van't Hoff).

\section{Ejercicios resueltos}\label{sec:ud2-resueltos}

Salvo que se indique otra cosa, usa las masas atómicas H = 1; C = 12; N = 14; O = 16; Na = 23; S = 32; Cl = 35,5; K = 39;
Ca = 40; y \(R = \qty{0.082}{\atm\litre\per\mol\per\kelvin}\).

\begin{example}
    \textbf{Ácido acético comercial.} Una botella de ácido acético, \ch{CH3COOH}, indica: «98\,\% en masa;
    \(d = \qty{1.05}{\gram\per\mL}\)». Calcula: a) su concentración molar; b) su concentración en \unit{\gram\per\litre};
    c) la molalidad de una disolución preparada mezclando \qty{15}{\mL} de este ácido con \qty{85}{\mL} de agua; d) el volumen
    de ácido comercial necesario para preparar \qty{500}{\mL} de disolución \qty{0.3}{\mol\per\litre}.

    \textbf{Solución.} \textbf{1. Datos}: riqueza 98\,\%, \(d = \qty{1.05}{\gram\per\mL}\), \(M = \qty{60}{\gram\per\mol}\).

    \textbf{2. Cálculos.} a) Partimos de \qty{100}{\gram} de disolución (98 g de ácido) y encadenamos factores:
    \[ c = \frac{\qty{98}{\gram}\ \text{ácido}}{\qty{100}{\gram}\ \text{dis.}} \cdot \frac{\qty{1}{\mol}}{\qty{60}{\gram}} \cdot
    \frac{\qty{1.05}{\gram}\ \text{dis.}}{\qty{1}{\mL}} \cdot \frac{\qty{1000}{\mL}}{\qty{1}{\litre}} = \qty{17.15}{\mol\per\litre} \]

    b) \(\gamma = \num{0.98} \cdot \qty{1050}{\gram\per\litre} = \qty{1029}{\gram\per\litre}\).

    c) Moles en los \qty{15}{\mL}: \(n = \num{17.15} \cdot \num{0.015} = \qty{0.257}{\mol}\). El disolvente son los \qty{85}{\gram}
    de agua añadidos (el agua que trae el ácido, \(15 \cdot \num{1.05} \cdot \num{0.02} = \qty{0.3}{\gram}\), es despreciable):
    \[ m = \frac{\qty{0.257}{\mol}}{\qty{0.085}{\kilo\gram}} = \qty{3.0}{\mol\per\kilo\gram} \]

    d) Moles necesarios: \(\num{0.3} \cdot \num{0.5} = \qty{0.15}{\mol}\); volumen de ácido comercial:
    \(V = \dfrac{\num{0.15}}{\num{17.15}} = \qty{8.75e-3}{\litre} = \qty{8.75}{\mL}\). Se toman con una pipeta, se vierten en un
    matraz aforado de \qty{500}{\mL} con algo de agua y se enrasa.
\end{example}

\begin{example}
    \textbf{Solubilidad del nitrato de sodio.} La solubilidad del nitrato de sodio a \qty{20}{\degreeCelsius} es de
    \qty{88}{\gram} por cada \qty{100}{\gram} de agua. a) ¿Cuántos gramos se pueden disolver como máximo en \qty{350}{\gram} de
    agua a esa temperatura? b) ¿Cuánto nitrato hay que añadir a \qty{0.6}{\kilo\gram} de agua para que la disolución quede
    saturada? c) ¿Será mayor o menor su solubilidad a \qty{60}{\degreeCelsius}?

    \textbf{Solución.} La solubilidad es una proporción: \qty{0.88}{\gram} de soluto por cada gramo de agua.

    a) \(\qty{350}{\gram}\ \text{agua} \cdot \dfrac{\qty{88}{\gram}}{\qty{100}{\gram}\ \text{agua}} = \qty{308}{\gram}\); si añadiéramos más,
    el exceso precipitaría.

    b) \(\qty{600}{\gram} \cdot \num{0.88} = \qty{528}{\gram}\) de nitrato de sodio.

    c) Como la de la mayoría de los sólidos, su solubilidad \textbf{aumenta} con la temperatura (a \qty{60}{\degreeCelsius} es de
    unos \qty{124}{\gram} por \qty{100}{\gram} de agua): las partículas se mueven más y la solvatación es más eficaz.
\end{example}

\begin{example}
    \textbf{Ley de Raoult.} A \qty{25}{\degreeCelsius}, la presión de vapor del agua es \qty{23.76}{\mmHg}, y la de una
    disolución acuosa de glucosa, \qty{21.40}{\mmHg}. a) ¿Cuál es la fracción molar de la glucosa? b) ¿Cuántos gramos de glucosa
    hay disueltos por cada \qty{100}{\gram} de agua?

    \textbf{Solución.} a) \(\Delta p = \num{23.76} - \num{21.40} = \qty{2.36}{\mmHg}\). Por la ley de Raoult:
    \[ \chi_s = \frac{\Delta p}{p^{\circ}} = \frac{\num{2.36}}{\num{23.76}} = \num{0.0993} \]
    La glucosa es el 9,9\,\% de los moles de la disolución, y el agua, el 90,1\,\%.

    b) En \qty{100}{\gram} de agua hay \(100/18 = \qty{5.56}{\mol}\). De \(\chi_s = \dfrac{n_s}{n_s + n_d}\):
    \(n_s = \dfrac{\chi_s \, n_d}{1 - \chi_s} = \dfrac{\num{0.0993} \cdot \num{5.56}}{\num{0.9007}} = \qty{0.613}{\mol}\), es decir,
    \(\num{0.613} \cdot 180 = \qty{110}{\gram}\) de glucosa.
\end{example}

\begin{example}
    \textbf{Mezcla de disoluciones.} Se mezclan \qty{150}{\mL} de una disolución de \ch{CaCl2} \qty{1.2}{\mol\per\litre} con
    \qty{250}{\mL} de otra \qty{2}{\mol\per\litre} del mismo soluto. a) ¿Cuál es la concentración de la mezcla? b) ¿Y si después
    se añaden \qty{100}{\mL} de agua? Supón volúmenes aditivos.

    \textbf{Solución.} a) Los moles de soluto se suman: \(n = \num{1.2} \cdot \num{0.15} + 2 \cdot \num{0.25} = \num{0.18} + \num{0.5} = \qty{0.68}{\mol}\),
    en \qty{0.4}{\litre}: \(c = \dfrac{\num{0.68}}{\num{0.4}} = \qty{1.7}{\mol\per\litre}\). La concentración queda entre las dos
    iniciales, más cerca de la de la disolución de la que se toma más volumen.

    b) El agua no aporta soluto: \(c = \dfrac{\num{0.68}}{\num{0.5}} = \qty{1.36}{\mol\per\litre}\).
\end{example}

\begin{example}
    \textbf{El anticongelante del coche.} El circuito de refrigeración de un coche contiene \qty{4}{\litre} de agua a los que se
    añaden \qty{0.5}{\litre} de etilenglicol, \ch{C2H6O2} (\(d = \qty{1.11}{\gram\per\mL}\)). Calcula las temperaturas de
    congelación y de ebullición del líquido.

    \textbf{Solución.} \(M(\ch{C2H6O2}) = \qty{62}{\gram\per\mol}\); \(m_{\text{etilenglicol}} = 500 \cdot \num{1.11} = \qty{555}{\gram}\);
    \(n = 555/62 = \qty{8.95}{\mol}\). Molalidad: \(m = \dfrac{\num{8.95}}{4} = \qty{2.24}{\mol\per\kilo\gram}\).
    \[ \Delta T_c = \num{1.86} \cdot \num{2.24} = \qty{4.16}{\degreeCelsius} \Rightarrow T_c' = \qty{-4.2}{\degreeCelsius}; \qquad
    \Delta T_e = \num{0.52} \cdot \num{2.24} = \qty{1.16}{\degreeCelsius} \Rightarrow T_e' = \qty{101.2}{\degreeCelsius} \]
    El anticongelante protege el motor de las heladas y también evita que el líquido hierva con facilidad.
\end{example}

%========================================================================================
\section{Ejercicios de aplicación}\label{sec:ud2-aplicacion}
%========================================================================================

\subsection{Porcentaje en masa, en volumen y concentración en masa}

\begin{exercise}
    \(\star\) Queremos preparar una disolución al 25\,\% en masa. a) Si tenemos \qty{40}{\gram} de soluto, ¿cuál será la masa de
    la disolución? b) ¿Cuánto soluto hay en \qty{60}{\gram} de disolución?
    \begin{solu}
        a) \(40 \cdot \dfrac{100}{25} = \textcolor{red}{\qty{160}{\gram}}\). b) \(60 \cdot \num{0.25} = \textcolor{red}{\qty{15}{\gram}}\).
    \end{solu}
\end{exercise}

\begin{exercise}
    \(\star\) Disolvemos \qty{40}{\gram} de \ch{KCl} en \qty{150}{\mL} de agua. ¿Cuál es el porcentaje en masa de la disolución?
    \begin{solu}
        \(\% = \dfrac{40}{190} \cdot 100 = \textcolor{red}{\qty{21.1}{\percent}}\).
    \end{solu}
\end{exercise}

\begin{exercise}
    \(\star\) Un medicamento es una disolución de \qty{200}{\gram\per\litre} y densidad \qty{1.15}{\gram\per\mL}. Calcula: a) su
    porcentaje en masa; b) la masa de soluto que hay en \qty{120}{\gram} de disolución; c) el volumen de disolución que contiene
    \qty{30}{\gram} de soluto.
    \begin{solu}
        a) En \qty{1}{\litre}: \qty{200}{\gram} de soluto en \qty{1150}{\gram}: \(\textcolor{red}{\qty{17.4}{\percent}}\).
        b) \(120 \cdot \num{0.174} = \textcolor{red}{\qty{20.9}{\gram}}\). c) \(30/200 = \qty{0.15}{\litre} = \textcolor{red}{\qty{150}{\mL}}\).
    \end{solu}
\end{exercise}

\begin{exercise}
    \(\star\star\) Una disolución de ibuprofeno tiene un 30\,\% en masa de principio activo y una densidad de
    \qty{1.25}{\gram\per\mL}. a) ¿Cuál es su concentración en \unit{\milli\gram\per\mL}? b) ¿Qué volumen contiene
    \qty{600}{\milli\gram} de ibuprofeno? c) ¿Cuánto ibuprofeno hay en \qty{4}{\mL}?
    \begin{solu}
        a) \(\num{0.30} \cdot 1250 = \textcolor{red}{\qty{375}{\milli\gram\per\mL}}\). b) \(600/375 = \textcolor{red}{\qty{1.6}{\mL}}\).
        c) \(4 \cdot 375 = \qty{1500}{\milli\gram} = \textcolor{red}{\qty{1.5}{\gram}}\).
    \end{solu}
\end{exercise}

\begin{exercise}
    \(\star\star\) La digoxina se usa para tratar arritmias cardíacas. A un paciente se le administran \qty{8}{\mL} de una
    disolución de digoxina de \qty{0.3}{\milli\gram} por cada \qty{5}{\mL}. ¿Cuál es la concentración en \unit{\gram\per\litre}?
    ¿Qué masa de digoxina ha recibido?
    \begin{solu}
        \(\gamma = \dfrac{\qty{0.3e-3}{\gram}}{\qty{5e-3}{\litre}} = \textcolor{red}{\qty{0.06}{\gram\per\litre}}\);
        \(m = 8 \cdot \dfrac{\num{0.3}}{5} = \textcolor{red}{\qty{0.48}{\milli\gram}}\).
    \end{solu}
\end{exercise}

\begin{exercise}
    \(\star\star\) Un vial de un anestésico local tiene una concentración del 2\,\% en masa. ¿Qué volumen hay que inyectar si la
    dosis prescrita es de \qty{120}{\milli\gram}? Supón que la densidad de la disolución es \qty{1}{\gram\per\mL}.
    \begin{solu}
        \(m_{\text{disolución}} = \dfrac{\num{0.120}}{\num{0.02}} = \qty{6}{\gram}\): \(\textcolor{red}{\qty{6}{\mL}}\).
    \end{solu}
\end{exercise}

\begin{exercise}
    \(\star\star\) ¿Cuántos gramos de glucosa hay en una bolsa de \qty{250}{\mL} de suero glucosado al 5\,\%? ¿Y de cloruro de
    sodio en una de \qty{500}{\mL} de suero fisiológico al 0,9\,\%? (Densidad de los sueros \(\approx \qty{1}{\gram\per\mL}\).)
    \begin{solu}
        \(250 \cdot \num{0.05} = \textcolor{red}{\qty{12.5}{\gram}}\) de glucosa; \(500 \cdot \num{0.009} = \textcolor{red}{\qty{4.5}{\gram}}\) de \ch{NaCl}.
    \end{solu}
\end{exercise}

\begin{exercise}
    \(\star\star\) La atropina se usa como antídoto en intoxicaciones por insecticidas. Se dispone de una disolución de
    \qty{4}{\gram} de atropina en \qty{100}{\mL}. ¿Cuál es su concentración en \unit{\gram\per\litre}? ¿Qué volumen hay que
    administrar para una dosis de \qty{200}{\milli\gram}?
    \begin{solu}
        \(\gamma = \textcolor{red}{\qty{40}{\gram\per\litre}}\) (\qty{40}{\milli\gram\per\mL}); \(V = 200/40 = \textcolor{red}{\qty{5}{\mL}}\).
    \end{solu}
\end{exercise}

\begin{exercise}
    \(\star\star\) A una niña de \qty{30}{\kilo\gram} hay que administrarle un medicamento a razón de \qty{20}{\milli\gram} por
    kilogramo de masa corporal. ¿Qué masa de medicamento le corresponde? Expresa la dosis como porcentaje en masa respecto a su
    masa corporal.
    \begin{solu}
        \(m = 30 \cdot 20 = \textcolor{red}{\qty{600}{\milli\gram}}\); \(\dfrac{\qty{20}{\milli\gram}}{\qty{1e6}{\milli\gram}} \cdot 100 =
        \textcolor{red}{\qty{2e-3}{\percent}}\).
    \end{solu}
\end{exercise}

\subsection{Concentración molar, molalidad y fracción molar}

\begin{exercise}
    \(\star\star\) Tenemos \qty{400}{\mL} de una disolución de \ch{HCl} \qty{0.4}{\mol\per\litre}. Determina: a) la masa de
    \ch{HCl} que contiene; b) el volumen de disolución que contiene \qty{365}{\milli\gram} de \ch{HCl}; c) su molalidad, si su
    densidad es \qty{1.02}{\gram\per\mL}.
    \begin{solu}
        a) \(n = \qty{0.16}{\mol}\): \(\textcolor{red}{\qty{5.84}{\gram}}\). b) \(\qty{0.01}{\mol}/\num{0.4} = \textcolor{red}{\qty{25}{\mL}}\).
        c) En \qty{1}{\litre}: \qty{1020}{\gram} de disolución, \(\num{0.4} \cdot \num{36.5} = \qty{14.6}{\gram}\) de \ch{HCl} y
        \qty{1005.4}{\gram} de agua: \(m = \textcolor{red}{\qty{0.398}{\mol\per\kilo\gram}}\).
    \end{solu}
\end{exercise}

\begin{exercise}
    \(\star\) Una mezcla contiene \qty{20}{\gram} de \ch{NaCl}, \qty{50}{\gram} de \ch{PbCl2} y \qty{10}{\gram} de \ch{NaOH}.
    Calcula la fracción molar de cada componente.
    \begin{solu}
        \(n(\ch{NaCl}) = \qty{0.342}{\mol}\); \(n(\ch{PbCl2}) = 50/\num{278.2} = \qty{0.180}{\mol}\); \(n(\ch{NaOH}) = \qty{0.25}{\mol}\); total \qty{0.772}{\mol}.
        \(\chi(\ch{NaCl}) = \textcolor{red}{\num{0.443}}\); \(\chi(\ch{PbCl2}) = \textcolor{red}{\num{0.233}}\); \(\chi(\ch{NaOH}) = \textcolor{red}{\num{0.324}}\).
    \end{solu}
\end{exercise}

\begin{exercise}
    \(\star\star\) En una disolución acuosa de \ch{KCl}, la fracción molar del soluto es \num{0.05}. Si se han usado
    \qty{400}{\gram} de agua, calcula: a) la masa de \ch{KCl}; b) la molalidad; c) la concentración molar, si la densidad es
    \qty{1.10}{\gram\per\mL}.
    \begin{solu}
        a) \(n_{\text{agua}} = \qty{22.2}{\mol}\); \(n_s = \dfrac{\num{0.05} \cdot \num{22.2}}{\num{0.95}} = \qty{1.17}{\mol}\):
        \(\textcolor{red}{\qty{87.3}{\gram}}\). b) \(\num{1.17}/\num{0.4} = \textcolor{red}{\qty{2.92}{\mol\per\kilo\gram}}\).
        c) \(V = \num{487.3}/\num{1.10} = \qty{443}{\mL}\): \(c = \textcolor{red}{\qty{2.64}{\mol\per\litre}}\).
    \end{solu}
\end{exercise}

\begin{exercise}
    \(\star\star\) Una botella de ácido sulfúrico comercial del 98\,\% en masa tiene una densidad de \qty{1.84}{\gram\per\mL}.
    Se toman \qty{8}{\mL} y se añade agua hasta tener \qty{1}{\litre} de disolución. Calcula la concentración molar: a) del ácido
    comercial; b) de la disolución preparada.
    \begin{solu}
        a) \(c = \dfrac{\num{0.98} \cdot 1840}{98} = \textcolor{red}{\qty{18.4}{\mol\per\litre}}\). b) \(\num{18.4} \cdot \num{0.008} = \textcolor{red}{\qty{0.147}{\mol\per\litre}}\).
    \end{solu}
\end{exercise}

\begin{exercise}
    \(\star\) Hay que preparar una disolución de curcumina, \ch{C21H20O6}, en etanol, de molalidad \qty{0.5}{\mol\per\kilo\gram}.
    ¿Qué masa de curcumina hay que añadir a \qty{200}{\gram} de etanol?
    \begin{solu}
        \(M = \qty{368}{\gram\per\mol}\); \(n = \num{0.5} \cdot \num{0.2} = \qty{0.1}{\mol}\): \(\textcolor{red}{\qty{36.8}{\gram}}\).
    \end{solu}
\end{exercise}

\begin{exercise}
    \(\star\star\) Disponemos de \qty{40}{\gram} de \ch{FeCl2}. ¿Qué masa de agua necesitamos para preparar: a) una disolución
    \qty{2}{\mol\per\kilo\gram}; b) una disolución en la que la fracción molar del agua sea \num{0.9}?
    \begin{solu}
        \(n = 40/\num{126.8} = \qty{0.315}{\mol}\). a) \(\num{0.315}/2 = \qty{0.158}{\kilo\gram} = \textcolor{red}{\qty{158}{\gram}}\).
        b) \(n_{\text{agua}} = 9 \, n_s = \qty{2.84}{\mol}\): \(\textcolor{red}{\qty{51.1}{\gram}}\).
    \end{solu}
\end{exercise}

\begin{exercise}
    \(\star\star\) En la farmacia me dan un frasco de \qty{1}{\litre} de una disolución de \qty{740}{\gram\per\litre} de un fármaco
    de masa molar \qty{185}{\gram\per\mol}. Determina: a) su concentración molar; b) la masa de fármaco del frasco; c) el volumen
    que debo tomar al día si la dosis es de \qty{370}{\milli\gram} dos veces al día.
    \begin{solu}
        a) \(740/185 = \textcolor{red}{\qty{4}{\mol\per\litre}}\). b) \(\textcolor{red}{\qty{740}{\gram}}\).
        c) \(\qty{740}{\milli\gram}\) diarios a \qty{0.74}{\gram\per\mL}: \(\textcolor{red}{\qty{1}{\mL}}\) al día.
    \end{solu}
\end{exercise}

\begin{exercise}
    \(\star\star\) En una fábrica de colonias se usa una disolución de etanol \qty{2.5}{\mol\per\litre}. a) ¿Qué masa de etanol
    hace falta para preparar \qty{400}{\mL}? b) ¿Qué volumen de disolución se obtiene con \qty{69}{\gram} de etanol?
    \begin{solu}
        a) \(n = \qty{1}{\mol}\): \(\textcolor{red}{\qty{46}{\gram}}\). b) \(n = 69/46 = \qty{1.5}{\mol}\): \(V = \num{1.5}/\num{2.5} = \textcolor{red}{\qty{0.6}{\litre}}\).
    \end{solu}
\end{exercise}

\begin{exercise}
    \(\star\star\) Un refresco en polvo, cuyo componente principal tiene una masa molar de \qty{342}{\gram\per\mol}, se prepara con
    una concentración máxima de \qty{0.2}{\mol\per\litre}. Cada sobre contiene \qty{34.2}{\gram}. Calcula: a) el volumen de
    bebida que se prepara con un sobre; b) la masa de polvo en \qty{1.5}{\litre} de bebida; c) qué fracción del sobre hay en un
    vaso de \qty{250}{\mL}.
    \begin{solu}
        a) \(n = \qty{0.1}{\mol}\): \(\textcolor{red}{\qty{0.5}{\litre}}\). b) \(\num{0.3} \cdot 342 = \textcolor{red}{\qty{102.6}{\gram}}\).
        c) \(\qty{0.05}{\mol} = \qty{17.1}{\gram}\): \(\textcolor{red}{\text{la mitad del sobre}}\).
    \end{solu}
\end{exercise}

\begin{exercise}
    \(\star\star\) Una bebida contiene un 12\,\% en masa de etanol, \ch{C2H6O}. Una persona de \qty{60}{\kilo\gram} bebe
    \qty{300}{\gram}, y el 15\,\% del alcohol pasa a la sangre. Calcula: a) la masa de etanol ingerida; b) la que pasa a la sangre;
    c) la concentración molar del etanol en la sangre (\qty{5}{\litre}); d) su molalidad en el cuerpo, considerando que toda la
    masa corporal actúa como disolvente.
    \begin{solu}
        a) \(\textcolor{red}{\qty{36}{\gram}}\). b) \(\textcolor{red}{\qty{5.4}{\gram}}\) (\qty{0.117}{\mol}).
        c) \(\num{0.117}/5 = \textcolor{red}{\qty{0.0235}{\mol\per\litre}}\). d) \(\num{0.117}/60 = \textcolor{red}{\qty{1.96e-3}{\mol\per\kilo\gram}}\).
    \end{solu}
\end{exercise}

\subsection{Preparación de disoluciones}

\begin{exercise}
    \(\star\) ¿Qué volumen de agua hay que añadir a \qty{250}{\mL} de \ch{HCl} \qty{0.3}{\mol\per\litre} para que su concentración
    sea \qty{0.05}{\mol\per\litre}? Supón volúmenes aditivos.
    \begin{solu}
        \(V_2 = \dfrac{\num{0.3} \cdot 250}{\num{0.05}} = \qty{1500}{\mL}\): hay que añadir \(\textcolor{red}{\qty{1.25}{\litre}}\).
    \end{solu}
\end{exercise}

\begin{exercise}
    \(\star\star\) Tenemos una disolución \qty{4}{\mol\per\kilo\gram} preparada con \qty{250}{\gram} de agua. a) ¿Cuánto soluto
    contiene si su masa molar es \qty{120}{\gram\per\mol}? b) ¿Cuánta agua hay que añadir para que la molalidad sea
    \qty{1.5}{\mol\per\kilo\gram}?
    \begin{solu}
        a) \(n = \qty{1}{\mol}\): \(\textcolor{red}{\qty{120}{\gram}}\). b) Agua total: \(1/\num{1.5} = \qty{0.667}{\kilo\gram}\);
        hay que añadir \(\textcolor{red}{\qty{417}{\gram}}\).
    \end{solu}
\end{exercise}

\begin{exercise}
    \(\star\) ¿Cómo prepararías \qty{250}{\mL} de alcohol de \ang{70} a partir de alcohol de \ang{96} y agua destilada?
    \begin{solu}
        \(V_1 = \dfrac{70 \cdot 250}{96} = \textcolor{red}{\qty{182}{\mL}}\) de alcohol de \ang{96}, \textcolor{red}{y se añade agua
        hasta completar \qty{250}{\mL}}.
    \end{solu}
\end{exercise}

\begin{exercise}
    \(\star\star\) ¿Cómo se prepara medio litro de un desinfectante de clorhexidina al 0,05\,\% (en volumen) a partir de
    clorhexidina al 20\,\%?
    \begin{solu}
        \(V_1 = \dfrac{\num{0.05} \cdot 500}{20} = \textcolor{red}{\qty{1.25}{\mL}}\) de la disolución al 20\,\%, \textcolor{red}{y se
        añade agua destilada hasta \qty{500}{\mL}}.
    \end{solu}
\end{exercise}

\begin{exercise}
    \(\star\star\) ¿Qué volumen de ácido nítrico del 65\,\% en masa y densidad \qty{1.40}{\gram\per\mL} hay que tomar para
    preparar \qty{250}{\mL} de disolución \qty{2}{\mol\per\litre}?
    \begin{solu}
        \(n = \qty{0.5}{\mol} = \qty{31.5}{\gram}\); \(m_{\text{comercial}} = \num{31.5}/\num{0.65} = \qty{48.5}{\gram}\):
        \(V = \textcolor{red}{\qty{34.6}{\mL}}\).
    \end{solu}
\end{exercise}

\begin{exercise}
    \(\star\star\) Se quiere preparar \qty{0.5}{\litre} de \ch{H2SO4} \qty{3}{\mol\per\litre} a partir de un ácido comercial de
    densidad \qty{1.84}{\gram\per\mL} y 96\,\% de riqueza en masa. ¿Qué volumen de ácido comercial hay que tomar? Explica cómo se
    prepararía.
    \begin{solu}
        \(n = \qty{1.5}{\mol} = \qty{147}{\gram}\); \(m_{\text{comercial}} = 147/\num{0.96} = \qty{153.1}{\gram}\):
        \(V = \textcolor{red}{\qty{83.2}{\mL}}\). \textcolor{blue}{Se vierte poco a poco, \textbf{sobre} agua, en un matraz aforado de
        \qty{500}{\mL}; se deja enfriar y se enrasa.}
    \end{solu}
\end{exercise}

\begin{exercise}
    \(\star\star\) Se desea preparar medio litro de \ch{HCl} \qty{0.4}{\mol\per\litre} a partir de un ácido comercial del 36\,\% en
    masa y densidad \qty{1.18}{\gram\per\mL}. ¿Qué volumen del comercial se necesita?
    \begin{solu}
        \(n = \qty{0.2}{\mol} = \qty{7.3}{\gram}\); \(m = \num{7.3}/\num{0.36} = \qty{20.3}{\gram}\): \(V = \textcolor{red}{\qty{17.2}{\mL}}\).
    \end{solu}
\end{exercise}

\begin{exercise}
    \(\star\star\) Se mezclan \qty{200}{\mL} de \ch{CaCl2} \qty{0.8}{\mol\per\litre} con \qty{300}{\mL} de \ch{CaCl2}
    \qty{1.5}{\mol\per\litre}. Suponiendo volúmenes aditivos, calcula la concentración: a) de la mezcla; b) si después se añaden
    \qty{250}{\mL} de agua.
    \begin{solu}
        a) \(n = \num{0.16} + \num{0.45} = \qty{0.61}{\mol}\): \(\num{0.61}/\num{0.5} = \textcolor{red}{\qty{1.22}{\mol\per\litre}}\).
        b) \(\num{0.61}/\num{0.75} = \textcolor{red}{\qty{0.81}{\mol\per\litre}}\).
    \end{solu}
\end{exercise}

\begin{exercise}
    \(\star\star\) Se mezclan \qty{150}{\mL} de \ch{HNO3} \qty{0.2}{\mol\per\litre} (disolución A) con \qty{100}{\mL} de \ch{HNO3}
    \qty{0.5}{\mol\per\litre} (B). Suponiendo volúmenes aditivos: a) calcula la concentración de la mezcla; b) ¿cuál sería si se
    añaden \qty{150}{\mL} de agua a la mezcla de A y B? c) ¿Y si, en lugar de eso, se evapora agua de la mezcla de A y B hasta
    reducir su volumen a la mitad?
    \begin{solu}
        \(n = \num{0.03} + \num{0.05} = \qty{0.08}{\mol}\). a) \(\num{0.08}/\num{0.25} = \textcolor{red}{\qty{0.32}{\mol\per\litre}}\).
        b) \(\num{0.08}/\num{0.4} = \textcolor{red}{\qty{0.2}{\mol\per\litre}}\). c) \(\num{0.08}/\num{0.125} = \textcolor{red}{\qty{0.64}{\mol\per\litre}}\).
    \end{solu}
\end{exercise}

\begin{exercise}
    \(\star\star\) Tenemos dos botellas: una de \ch{NaOH} \qty{0.4}{\mol\per\litre} y otra de \ch{KOH} \qty{0.4}{\mol\per\litre}.
    Echamos el mismo volumen de cada una en dos vasos. a) ¿Qué vaso tiene más moles de soluto? b) ¿Y más masa de soluto?
    \begin{solu}
        a) \textcolor{red}{Los dos tienen los mismos moles} (misma concentración y mismo volumen). b) \textcolor{red}{El de \ch{KOH}},
        porque su masa molar (\qty{56}{\gram\per\mol}) es mayor que la del \ch{NaOH} (\qty{40}{\gram\per\mol}).
    \end{solu}
\end{exercise}

\begin{exercise}
    \(\star\star\) Un recipiente contiene \qty{3}{\litre} de \ch{KCl} \qty{2}{\mol\per\litre} y otro, \qty{2}{\litre} de \ch{KCl}
    \qty{4}{\mol\per\litre}. a) ¿Cuál tiene más masa de \ch{KCl}? b) ¿Cuál es la concentración si se mezclan? (Volúmenes aditivos.)
    \begin{solu}
        a) \(6 \cdot \num{74.6} = \qty{448}{\gram}\) frente a \(8 \cdot \num{74.6} = \qty{597}{\gram}\): \textcolor{red}{el segundo}.
        b) \(14/5 = \textcolor{red}{\qty{2.8}{\mol\per\litre}}\).
    \end{solu}
\end{exercise}

\begin{exercise}
    \(\star\star\) El agua oxigenada comercial es una disolución de peróxido de hidrógeno al 3\,\% en volumen. Para un enjuague bucal
    se necesita una disolución al 1,5\,\%. Tenemos un envase de \qty{500}{\mL} de agua oxigenada y uno de \qty{1}{\litre} de agua
    destilada. Suponiendo volúmenes aditivos: a) ¿cuánta agua oxigenada y cuánta agua hacen falta para \qty{200}{\mL} de enjuague?
    b) ¿Y para \qty{750}{\mL}? c) Después de preparar las dos, ¿cuánto queda en cada envase?
    \begin{solu}
        Se diluye a la mitad: a) \(\textcolor{red}{\qty{100}{\mL}}\) de cada. b) \(\textcolor{red}{\qty{375}{\mL}}\) de cada.
        c) Agua oxigenada: \(500 - 475 = \textcolor{red}{\qty{25}{\mL}}\); agua: \(1000 - 475 = \textcolor{red}{\qty{525}{\mL}}\).
    \end{solu}
\end{exercise}

\begin{exercise}
    \(\star\star\) Se quiere preparar \qty{1.5}{\litre} de \ch{HCl} \qty{0.2}{\mol\per\litre} a partir de un ácido comercial del
    37\,\% en masa y densidad \qty{1.19}{\gram\per\cubic\cm}. Calcula: a) la concentración molar del ácido comercial; b) el volumen
    de comercial necesario.
    \begin{solu}
        a) \(\dfrac{\num{0.37} \cdot 1190}{\num{36.5}} = \textcolor{red}{\qty{12.1}{\mol\per\litre}}\). b) \(\dfrac{\num{0.3}}{\num{12.06}} = \textcolor{red}{\qty{24.9}{\mL}}\).
    \end{solu}
\end{exercise}

\begin{exercise}
    \(\star\star\) Se toman \qty{15}{\mL} de \ch{H2SO4} comercial del 98\,\% en masa y densidad \qty{1.84}{\gram\per\mL}, y se
    añaden a un matraz de \qty{500}{\mL} con agua, que después se enrasa. Calcula la concentración molar: a) del ácido comercial;
    b) de la disolución preparada.
    \begin{solu}
        a) \(\textcolor{red}{\qty{18.4}{\mol\per\litre}}\). b) \(\dfrac{\num{18.4} \cdot \num{0.015}}{\num{0.5}} = \textcolor{red}{\qty{0.552}{\mol\per\litre}}\).
    \end{solu}
\end{exercise}

\subsection{Solubilidad}

\begin{exercise}
    \(\star\) ¿Qué es la solubilidad? ¿En qué unidades se expresa? ¿Qué factores la modifican?
    \begin{solu}
        \textcolor{red}{Es la máxima cantidad de soluto que se disuelve en una cantidad dada de disolvente, a una temperatura y una
        presión dadas}. Se expresa en g de soluto/\qty{100}{\gram} de disolvente o en \unit{\gram\per\litre}. Depende de la
        naturaleza del soluto y del disolvente, de la temperatura y, en los gases, de la presión.
    \end{solu}
\end{exercise}

\begin{exercise}
    \(\star\star\) Explica razonadamente qué café quedará más dulce: a) un café frío al que se añade un terrón de azúcar y no se
    remueve; b) un café caliente con una cucharada de azúcar en polvo, removiendo.
    \begin{solu}
        \textcolor{red}{El b}: el azúcar en polvo tiene más superficie de contacto, al agitar se renueva el disolvente junto al
        soluto y el calor acelera la disolución y aumenta la solubilidad. En a, buena parte del azúcar puede quedar sin disolver en
        el fondo.
    \end{solu}
\end{exercise}

\begin{exercise}
    \(\star\star\) Con las curvas de solubilidad de la figura~\ref{fig:ud2-curvas}: a) ¿Qué sustancia es la más soluble a
    \qty{10}{\degreeCelsius}? ¿Y a \qty{60}{\degreeCelsius}? b) ¿Cuál es la solubilidad del \ch{KNO3} a 20 y a
    \qty{50}{\degreeCelsius}? ¿Por qué se dice que es un gran cambio? c) ¿Qué sustancias pueden disolverse al menos
    \qty{40}{\gram} por \qty{100}{\gram} de agua a \qty{20}{\degreeCelsius}? ¿Y a \qty{60}{\degreeCelsius}?
    \begin{solu}
        a) \textcolor{red}{A \qty{10}{\degreeCelsius}, el \ch{CaCl2}}; \textcolor{red}{a \qty{60}{\degreeCelsius}, el \ch{KNO3}}
        (unos \qty{110}{\gram}). b) \textcolor{red}{Unos \qty{32}{\gram} y \qty{85}{\gram}}: casi se triplica en solo
        \qty{30}{\degreeCelsius}. c) \textcolor{red}{A \qty{20}{\degreeCelsius}, solo el \ch{CaCl2}; a \qty{60}{\degreeCelsius},
        \ch{KNO3}, \ch{CaCl2}, \ch{K2Cr2O7} y \ch{KCl}}.
    \end{solu}
\end{exercise}

\begin{exercise}
    \(\star\star\) ¿Cómo podrías disminuir la solubilidad del azúcar en agua? ¿Para qué puede servir?
    \begin{solu}
        \textcolor{red}{Enfriando la disolución}, ya que su solubilidad aumenta con la temperatura. Sirve para \textbf{cristalizar}:
        al enfriar una disolución saturada en caliente, el exceso de soluto se separa en forma de cristales puros (así se obtiene,
        por ejemplo, el azúcar de las refinerías).
    \end{solu}
\end{exercise}

\begin{exercise}
    \(\star\star\) Infórmate sobre el «mal de altura» o hipoxia. Explícalo con la ley de Henry.
    \begin{solu}
        \textcolor{red}{En altura la presión atmosférica es menor, y con ella la presión parcial del oxígeno}. Por la ley de Henry,
        se disuelve menos oxígeno en la sangre y llega menos a los tejidos: aparecen dolor de cabeza, cansancio y náuseas.
    \end{solu}
\end{exercise}

\begin{exercise}
    \(\star\star\) Los buceadores que descienden a gran profundidad no respiran aire, sino mezclas de oxígeno con helio. ¿Por qué?
    \begin{solu}
        \textcolor{red}{A gran profundidad la presión es muy alta y, por la ley de Henry, se disolvería mucho nitrógeno en la
        sangre}, que produce narcosis y, al subir, burbujas (enfermedad descompresiva). \textcolor{red}{El helio es mucho menos
        soluble en la sangre}, así que el riesgo es menor.
    \end{solu}
\end{exercise}

\begin{exercise}
    \(\star\star\) Una lata de \qty{330}{\mL} de refresco se envasa con una presión de \ch{CO2} de \qty{4}{\atm}. La constante de
    Henry del \ch{CO2} en agua a esa temperatura es \qty{3.4e-2}{\mol\per\litre\per\atm}. a) ¿Qué masa de \ch{CO2} hay disuelta?
    b) ¿Qué ocurre al abrirla, si en el aire la presión parcial del \ch{CO2} es de unos \qty{4e-4}{\atm}?
    \begin{solu}
        a) \(c = \num{3.4e-2} \cdot 4 = \qty{0.136}{\mol\per\litre}\); \(n = \qty{0.0449}{\mol}\): \(\textcolor{red}{\qty{1.97}{\gram}}\).
        b) \textcolor{red}{La solubilidad pasa a ser unas \num{10000} veces menor}: casi todo el \ch{CO2} acaba escapando (el refresco
        «se queda sin gas»).
    \end{solu}
\end{exercise}

\subsection{Propiedades coligativas}

\begin{exercise}
    \(\star\star\) Calcula la presión de vapor a \qty{25}{\degreeCelsius} de una disolución de \qty{80}{\gram} de glucosa en
    \qty{300}{\gram} de agua (\(p^{\circ} = \qty{23.8}{\mmHg}\)).
    \begin{solu}
        \(n_s = \qty{0.444}{\mol}\); \(n_d = \qty{16.67}{\mol}\); \(\chi_s = \num{0.026}\): \(\Delta p = \qty{0.62}{\mmHg}\);
        \(p = \textcolor{red}{\qty{23.2}{\mmHg}}\).
    \end{solu}
\end{exercise}

\begin{exercise}
    \(\star\star\) Se disuelven \qty{20}{\gram} de un soluto no volátil en \qty{250}{\gram} de tolueno, \ch{C6H5CH3}. La presión de
    vapor de la disolución es \qty{24.5}{\mmHg} y la del tolueno puro, \qty{28.12}{\mmHg}. Calcula la masa molar del soluto.
    \begin{solu}
        \(\chi_s = \dfrac{\num{3.62}}{\num{28.12}} = \num{0.129}\); \(n_d = 250/92 = \qty{2.72}{\mol}\);
        \(n_s = \dfrac{\num{0.129} \cdot \num{2.72}}{\num{0.871}} = \qty{0.402}{\mol}\): \(M = \textcolor{red}{\qty{49.8}{\gram\per\mol}}\).
    \end{solu}
\end{exercise}

\begin{exercise}
    \(\star\star\) Se disuelven \qty{25}{\gram} de un sólido no volátil en \qty{200}{\gram} de benceno, \ch{C6H6}. A
    \qty{25}{\degreeCelsius}, la presión de vapor de la disolución es \qty{84.5}{\mmHg} y la del benceno, \qty{95.1}{\mmHg}.
    Calcula: a) la masa molar del soluto; b) la molalidad de la disolución.
    \begin{solu}
        \(\chi_s = \num{10.6}/\num{95.1} = \num{0.111}\); \(n_d = \qty{2.56}{\mol}\); \(n_s = \qty{0.322}{\mol}\).
        a) \(M = \textcolor{red}{\qty{77.7}{\gram\per\mol}}\). b) \(\num{0.322}/\num{0.2} = \textcolor{red}{\qty{1.61}{\mol\per\kilo\gram}}\).
    \end{solu}
\end{exercise}

\begin{exercise}
    \(\star\star\) Al disolver \qty{1.5}{\gram} de resorcina en \qty{50}{\gram} de agua, la temperatura de ebullición aumenta
    \qty{0.142}{\degreeCelsius}. ¿Cuál es la masa molar de la resorcina?
    \begin{solu}
        \(m = \num{0.142}/\num{0.52} = \qty{0.273}{\mol\per\kilo\gram}\); \(n = \qty{0.01365}{\mol}\): \(M = \textcolor{red}{\qty{110}{\gram\per\mol}}\).
    \end{solu}
\end{exercise}

\begin{exercise}
    \(\star\) Una disolución acuosa de un soluto no volátil tiene una molalidad de \qty{2}{\mol\per\kilo\gram}. Calcula sus puntos
    de ebullición y de congelación.
    \begin{solu}
        \(T_e' = 100 + \num{0.52} \cdot 2 = \textcolor{red}{\qty{101.04}{\degreeCelsius}}\); \(T_c' = 0 - \num{1.86} \cdot 2 = \textcolor{red}{\qty{-3.72}{\degreeCelsius}}\).
    \end{solu}
\end{exercise}

\begin{exercise}
    \(\star\star\) Una disolución de \qty{3}{\gram} de un soluto en \qty{80}{\gram} de alcanfor congela a \qty{166}{\degreeCelsius}.
    Calcula la masa molar del soluto (consulta la tabla~\ref{tab:ud2-constantes}).
    \begin{solu}
        \(\Delta T_c = 178 - 166 = \qty{12}{\degreeCelsius}\); \(m = 12/40 = \qty{0.3}{\mol\per\kilo\gram}\); \(n = \qty{0.024}{\mol}\):
        \(M = \textcolor{red}{\qty{125}{\gram\per\mol}}\).
    \end{solu}
\end{exercise}

\begin{exercise}
    \(\star\) Calcula la presión osmótica de una disolución de \qty{20}{\gram} de glucosa en \qty{250}{\mL} de disolución a
    \qty{25}{\degreeCelsius}.
    \begin{solu}
        \(c = \dfrac{20/180}{\num{0.25}} = \qty{0.444}{\mol\per\litre}\): \(\pi = \num{0.444} \cdot \num{0.082} \cdot 298 = \textcolor{red}{\qty{10.9}{\atm}}\).
    \end{solu}
\end{exercise}

\begin{exercise}
    \(\star\star\) Calcula la presión osmótica a \qty{25}{\degreeCelsius} de una disolución de sacarosa, \ch{C12H22O11}, del 30\,\% en
    masa y densidad \qty{1.13}{\gram\per\mL}.
    \begin{solu}
        \(c = \dfrac{\num{0.30} \cdot 1130}{342} = \qty{0.991}{\mol\per\litre}\): \(\pi = \textcolor{red}{\qty{24.2}{\atm}}\).
    \end{solu}
\end{exercise}

\begin{exercise}
    \(\star\star\) Calcula la masa molar de un soluto sabiendo que una disolución de \qty{4}{\gram} de soluto en \qty{1.5}{\litre}
    de disolución tiene, a \qty{25}{\degreeCelsius}, una presión osmótica de \qty{0.80}{\atm}.
    \begin{solu}
        \(M = \dfrac{m R T}{\pi V} = \dfrac{4 \cdot \num{0.082} \cdot 298}{\num{0.80} \cdot \num{1.5}} = \textcolor{red}{\qty{81.5}{\gram\per\mol}}\).
    \end{solu}
\end{exercise}

\begin{exercise}
    \(\star\star\) ¿A qué temperatura la presión osmótica de una disolución \qty{0.05}{\mol\per\litre} de urea es de
    \qty{950}{\mmHg}?
    \begin{solu}
        \(\pi = \qty{1.25}{\atm}\); \(T = \dfrac{\pi}{c R} = \dfrac{\num{1.25}}{\num{0.05} \cdot \num{0.082}} = \textcolor{red}{\qty{305}{\kelvin}}\) (\qty{32}{\degreeCelsius}).
    \end{solu}
\end{exercise}

\begin{exercise}
    \(\star\) ¿Qué es la ósmosis? ¿Qué diferencia hay entre ósmosis y ósmosis inversa?
    \begin{solu}
        \textcolor{red}{La ósmosis es el paso espontáneo de disolvente a través de una membrana semipermeable, de la disolución más
        diluida a la más concentrada}. \textcolor{red}{En la ósmosis inversa se aplica a la disolución concentrada una presión mayor
        que la osmótica}, y el disolvente pasa en sentido contrario (así funcionan las desaladoras).
    \end{solu}
\end{exercise}

\begin{exercise}
    \(\star\) ¿Por qué el pescado en salazón se conserva mucho tiempo y en la nevera solo unos días?
    \begin{solu}
        \textcolor{red}{La sal crea un medio hipertónico}: los microorganismos pierden agua por ósmosis y no pueden crecer. En la
        nevera el frío solo frena su crecimiento.
    \end{solu}
\end{exercise}

\begin{exercise}
    \(\star\star\) ¿Por qué, para quitar el exceso de sal del bacalao o del jamón, se dejan en agua?
    \begin{solu}
        \textcolor{red}{El agua es un medio hipotónico respecto al alimento}: la sal pasa al agua (y el agua al alimento) hasta que
        las concentraciones tienden a igualarse. Cambiando el agua varias veces se elimina más sal.
    \end{solu}
\end{exercise}

\begin{exercise}
    \(\star\star\) La concentración salina de la sangre de un paciente deshidratado es más alta de lo normal. ¿Qué tipo de suero
    le pondrías para reducirla, isotónico, hipotónico o hipertónico?
    \begin{solu}
        \textcolor{red}{Un suero hipotónico} (con menos sales que la sangre), administrado con cuidado: aporta agua y reduce la
        concentración. Uno hipertónico la aumentaría aún más.
    \end{solu}
\end{exercise}

\begin{exercise}
    \(\star\star\) A \qty{37}{\degreeCelsius}, el plasma sanguíneo tiene la misma concentración de partículas que el interior de los
    glóbulos rojos, unos \qty{0.3}{\mol\per\litre}. Explica qué les ocurre a los glóbulos rojos si se introducen en: a) un suero de
    \qty{1}{\mol\per\litre}; b) un suero de \qty{0.1}{\mol\per\litre}; c) un suero de \qty{0.3}{\mol\per\litre}.
    \begin{solu}
        a) \textcolor{red}{Medio hipertónico: pierden agua y se arrugan}. b) \textcolor{red}{Medio hipotónico: ganan agua, se hinchan
        y pueden romperse (hemólisis)}. c) \textcolor{red}{Medio isotónico: no cambian}.
    \end{solu}
\end{exercise}

%========================================================================================
\section{Ejercicios de ampliación}\label{sec:ud2-ampliacion}
%========================================================================================

\subsection{¿Quieres ser enfermero o enfermera por un día?}

Una dosis expresada en \unit{\mL\per\kilo\gram\per\minuto} indica cuántos mililitros hay que administrar por cada kilogramo
de masa corporal y por cada minuto, es decir, \(\dfrac{\unit{\mL}}{\unit{\kilo\gram} \cdot \unit{\minuto}}\).

\begin{exercise}
    \(\star\star\) Se dispone de \qty{250}{\mL} de un jarabe con una concentración de \qty{300}{\milli\gram} por cada
    \qty{5}{\mL}. Un paciente debe tomar \qty{2.4}{\gram} al día en cuatro dosis. a) ¿Cuál es la concentración en
    \unit{\gram\per\litre}? b) ¿Cuántos mililitros debe tomar en cada dosis?
    \begin{solu}
        a) \(\textcolor{red}{\qty{60}{\gram\per\litre}}\). b) \qty{600}{\milli\gram} por dosis a \qty{60}{\milli\gram\per\mL}:
        \(\textcolor{red}{\qty{10}{\mL}}\).
    \end{solu}
\end{exercise}

\begin{exercise}
    \(\star\star\) Un paciente debe recibir \qty{1.2}{\gram} diarios de un analgésico, disuelto al 20\,\% en masa, en perfusión
    continua. ¿Cuántos miligramos de disolución por minuto hay que administrarle?
    \begin{solu}
        \(m_{\text{disolución}} = \num{1.2}/\num{0.2} = \qty{6}{\gram} = \qty{6000}{\milli\gram}\) en \qty{1440}{\minuto}:
        \(\textcolor{red}{\qty{4.17}{\milli\gram\per\minuto}}\).
    \end{solu}
\end{exercise}

\begin{exercise}
    \(\star\star\star\) Una persona de \qty{80}{\kilo\gram} recibe una disolución de un medicamento al 10\,\%, de modo que le llegan
    \qty{0.15}{\milli\gram} de medicamento por kilogramo y minuto. ¿Cuál es el ritmo de perfusión de la disolución en
    \unit{\milli\gram\per\hour}?
    \begin{solu}
        Medicamento: \(80 \cdot \num{0.15} \cdot 60 = \qty{720}{\milli\gram\per\hour}\); disolución: \(720/\num{0.1} = \textcolor{red}{\qty{7200}{\milli\gram\per\hour}}\).
    \end{solu}
\end{exercise}

\begin{exercise}
    \(\star\star\star\) Un niño de \qty{25}{\kilo\gram} recibe una disolución de un medicamento al 10\,\% a razón de
    \qty{0.2}{\milli\gram} de medicamento por kilogramo y minuto. ¿Cuál es el ritmo de perfusión en \unit{\milli\gram\per\hour}?
    \begin{solu}
        \(\dfrac{25 \cdot \num{0.2} \cdot 60}{\num{0.1}} = \textcolor{red}{\qty{3000}{\milli\gram\per\hour}}\).
    \end{solu}
\end{exercise}

\begin{exercise}
    \(\star\star\) En 24 horas hay que pasar a un paciente \qty{1500}{\cubic\cm} de suero glucosado junto con \qty{50}{\cubic\cm}
    de \ch{KCl} \qty{2}{\mol\per\litre}. a) ¿Cuántos sueros se usarán si los envases son de \qty{500}{\cubic\cm}? b) ¿Cuántos
    mililitros de \ch{KCl} hay que poner en cada suero? c) ¿A qué ritmo de perfusión, en \unit{\mL\per\hour}, hay que mantenerlos?
    \begin{solu}
        a) \(\textcolor{red}{3}\) envases. b) \(50/3 = \textcolor{red}{\qty{16.7}{\mL}}\). c) \(1550/24 = \textcolor{red}{\qty{64.6}{\mL\per\hour}}\).
    \end{solu}
\end{exercise}

\begin{exercise}
    \(\star\star\) A Lucía, que pesa \qty{18}{\kilo\gram}, hay que darle \qty{25}{\milli\gram} por kilogramo y día de un
    medicamento, repartidos en dos dosis. a) ¿Cuántos mililitros de un jarabe de \qty{5}{\milli\gram\per\mL} hay que darle en
    cada dosis? b) Expresa la dosis en \unit{\mL\per\kilo\gram} al día.
    \begin{solu}
        a) \(18 \cdot 25/2 = \qty{225}{\milli\gram}\) por dosis: \(\textcolor{red}{\qty{45}{\mL}}\). b) \(25/5 = \textcolor{red}{\qty{5}{\mL\per\kilo\gram}}\) al día.
    \end{solu}
\end{exercise}

\begin{exercise}
    \(\star\star\) Para una otitis, a un niño de \qty{14}{\kilo\gram} se le receta amoxicilina a \qty{80}{\milli\gram} por
    kilogramo y día, en tres tomas. ¿Cuántos mililitros de un jarabe de \qty{250}{\milli\gram} por cada \qty{5}{\mL} hay que darle
    en cada toma?
    \begin{solu}
        \(14 \cdot 80/3 = \qty{373}{\milli\gram}\) por toma, a \qty{50}{\milli\gram\per\mL}: \(\textcolor{red}{\qty{7.5}{\mL}}\).
    \end{solu}
\end{exercise}

\begin{exercise}
    \(\star\star\star\) Los preparados homeopáticos se obtienen por diluciones sucesivas: en uno «1 CH» se toma una parte de la
    disolución inicial y se diluye en 99 partes de disolvente; en uno «2 CH» se repite la operación con el anterior, y así
    sucesivamente. Partiendo de una disolución \qty{1}{\mol\per\litre}: a) ¿cuál es la concentración de un preparado 2 CH? ¿Y de
    uno 30 CH? b) ¿Cuántas moléculas de principio activo hay en \qty{100}{\mL} de un preparado 12 CH? ¿Qué conclusión sacas?
    \begin{solu}
        a) Cada CH divide la concentración entre 100: \(\textcolor{red}{\qty{e-4}{\mol\per\litre}}\) y \(\textcolor{red}{\qty{e-60}{\mol\per\litre}}\).
        b) \(\num{e-24} \cdot \num{0.1} \cdot \num{6.022e23} = \textcolor{red}{\num{0.06}}\): \textcolor{red}{lo más probable es que no
        haya ni una sola molécula} de principio activo; el preparado es solo disolvente.
    \end{solu}
\end{exercise}

\subsection{Propiedades coligativas}

\begin{exercise}
    \(\star\star\star\) Al disolver \qty{30}{\gram} de una sustancia desconocida en \qty{150}{\gram} de agua, la presión de vapor
    pasa de \qty{23.76}{\mmHg} a \qty{23.36}{\mmHg}. Calcula: a) la variación de la presión de vapor; b) la masa molar de la
    sustancia.
    \begin{solu}
        a) \(\textcolor{red}{\qty{0.40}{\mmHg}}\). b) \(\chi_s = \num{0.40}/\num{23.76} = \num{0.0168}\); \(n_d = \qty{8.33}{\mol}\);
        \(n_s = \qty{0.143}{\mol}\): \(M = \textcolor{red}{\qty{210}{\gram\per\mol}}\).
    \end{solu}
\end{exercise}

\begin{exercise}
    \(\star\star\) Una disolución contiene un 40\,\% en masa de un soluto de masa molar \qty{902}{\gram\per\mol} en cloroformo.
    ¿Cuánto aumenta su temperatura de ebullición? (Consulta la tabla~\ref{tab:ud2-constantes}.)
    \begin{solu}
        En \qty{100}{\gram}: \(n_s = 40/902 = \qty{0.0443}{\mol}\) en \qty{0.060}{\kilo\gram}: \(m = \qty{0.739}{\mol\per\kilo\gram}\);
        \(\Delta T_e = \num{3.63} \cdot \num{0.739} = \textcolor{red}{\qty{2.7}{\degreeCelsius}}\).
    \end{solu}
\end{exercise}

\begin{exercise}
    \(\star\star\) El radiador de un coche contiene \qty{5}{\litre} de agua y \qty{400}{\mL} de etilenglicol, \ch{C2H6O2}
    (\(d = \qty{1.11}{\gram\per\mL}\)). ¿Cuál es la temperatura más baja que puede soportar sin que se congele?
    \begin{solu}
        \(n = \dfrac{400 \cdot \num{1.11}}{62} = \qty{7.16}{\mol}\); \(m = \qty{1.43}{\mol\per\kilo\gram}\); \(\Delta T_c = \qty{2.66}{\degreeCelsius}\):
        \(\textcolor{red}{\qty{-2.7}{\degreeCelsius}}\). \textcolor{blue}{Con tan poco anticongelante, el coche está poco protegido.}
    \end{solu}
\end{exercise}

\begin{exercise}
    \(\star\star\) Calcula la presión osmótica a \qty{37}{\degreeCelsius} de una disolución de ácido ascórbico (vitamina C),
    \ch{C6H8O6}, que contiene \qty{150}{\gram\per\litre}.
    \begin{solu}
        \(c = 150/176 = \qty{0.852}{\mol\per\litre}\): \(\pi = \num{0.852} \cdot \num{0.082} \cdot 310 = \textcolor{red}{\qty{21.7}{\atm}}\).
    \end{solu}
\end{exercise}

\begin{exercise}
    \(\star\star\star\) El agua de mar contiene, aproximadamente, un 3,5\,\% en masa de sales, que supondremos que son todas
    \ch{NaCl}. Teniendo en cuenta que cada unidad de \ch{NaCl} da dos iones (\(i = 2\)), estima la temperatura de congelación del
    agua de mar. ¿Por qué en el mar el hielo se forma más tarde que en los lagos?
    \begin{solu}
        En \qty{1000}{\gram} de agua de mar: \(n = 35/\num{58.5} = \qty{0.598}{\mol}\) de \ch{NaCl} en \qty{0.965}{\kilo\gram} de agua:
        \(m = \qty{0.620}{\mol\per\kilo\gram}\). \(\Delta T_c = i K_c m = 2 \cdot \num{1.86} \cdot \num{0.620} = \qty{2.3}{\degreeCelsius}\):
        \(\textcolor{red}{T_c' \approx \qty{-2.3}{\degreeCelsius}}\) (el valor real es algo mayor, unos \qty{-1.9}{\degreeCelsius}).
        \textcolor{red}{Las sales disueltas bajan su punto de congelación}.
    \end{solu}
\end{exercise}

\begin{exercise}
    \(\star\star\star\) El suero fisiológico es una disolución de \ch{NaCl} al 0,9\,\% (\qty{9}{\gram\per\litre}), isotónica con la
    sangre. a) Calcula su presión osmótica a \qty{37}{\degreeCelsius}, teniendo en cuenta que \(i = 2\). b) ¿Qué concentración de
    glucosa, en \unit{\gram\per\litre}, debe tener un suero glucosado para ser también isotónico? Compárala con la de los sueros
    glucosados del 5\,\%.
    \begin{solu}
        a) \(c = 9/\num{58.5} = \qty{0.154}{\mol\per\litre}\); partículas: \(2c = \qty{0.308}{\mol\per\litre}\);
        \(\pi = \num{0.308} \cdot \num{0.082} \cdot 310 = \textcolor{red}{\qty{7.8}{\atm}}\). b) La glucosa no se disocia: necesita
        \(\qty{0.308}{\mol\per\litre} \cdot 180 = \textcolor{red}{\qty{55}{\gram\per\litre}}\), es decir, \textcolor{red}{un 5,5\,\%},
        muy próximo al 5\,\% de los sueros glucosados.
    \end{solu}
\end{exercise}
